% Part: first-order-logic
% Chapter: lindstrom
% Section: lindstrom-proof

\documentclass[../../../include/open-logic-section]{subfiles}

\begin{document}

\olfileid{mod}{lin}{prf}

\olsection{లిండ్‌స్ట్రోమ్ సిద్ధాంతం}

\begin{lem}
\ollabel{lem:lindstrom}
సాధారణ నైరూప్య తర్కంలో $!E\in L(\Lang{L})$
అనుకుందాం; $\Lang{L}$ పరిమిత భాష. ఇంకా, ఏ రెండు
\tetoken{నిర్మాణాలు}{!!{structure}s} $\Struct{M}$, $\Struct{N}$లకు
$\Struct{M}\equiv_n\Struct{N}$, $\Struct{M}\models_L!E$ అయితే
$\Struct{N}\models_L!E$ అయ్యే $n\in\Nat$ ఉందనుకుందాం. అప్పుడు $!E$
మొదటిస్థాయి \tetoken{వాక్యానికి}{!!{sentence}} తుల్యం; అంటే,
మొదటిస్థాయి $!D$ ఉండి $\Mod(L){!E}=\Mod(L){!D}$.
\sourcecorrection{OLTEMODLIN-008}{ఈ లెమ్మా తరువాత నిర్మించే
మొదటిస్థాయి $!D$ను $L$-వాక్యంగా తీసుకుని, తుల్యతను సాధారణ
తర్కపు సంతృప్తితో పోల్చాలంటే సాధారణత్వం అవసరం. ముందరి విస్తరణ
ధర్మం ద్వారా మొదటిస్థాయి తర్కం ప్రతి సాధారణ తర్కంలో ఇమిడుతుందని
చెప్పినందున, లెమ్మా పరికల్పనలో “సాధారణ”ను స్పష్టం చేశాం.}
\end{lem}

\begin{proof}
$!E$ విలువపై ఏ రెండు $n$-తుల్య \tetoken{నిర్మాణాలు}{!!{structure}s}
$\Struct{M}$, $\Struct{N}$
ఏకీభవించేలా $n$ను ఎంచుకోండి. \olref[bas][pis]{prop:qr-finite}ను
గుర్తుచేసుకుంటే, పరిమాణక ర్యాంకు $n$కు మించని మొదటిస్థాయి
\tetoken{వాక్యాల}{!!{sentence}s} తార్కిక తుల్యతా-వర్గాలు, పరిమిత
\tetoken{భాష}{!!{language}}లో,
పరిమిత సంఖ్యలోనే ఉంటాయి. ప్రతి స్థిర
\tetoken{నిర్మాణం}{!!{structure}}~$\Struct{M}$కు, $\Struct{M}$లో
సత్యమైన $!E$, $\QuantRank{!A}\le n$ గల మొదటిస్థాయి
\tetoken{వాక్యాల}{!!{sentence}s} తుల్యతా-వర్గాల నుంచి ఒక్కో
ప్రతినిధిని తీసుకొని, వాటి సంయోగాన్ని $!D_{\Struct{M}}$గా ఉంచండి.
ఈ సంయోగం పరిమితమే; అప్పుడు $\Struct{N}\models!D_{\Struct{M}}$
అవుతుంది అప్పుడే $\Struct{N}\equiv_n\Struct{M}$.
\sourcecorrection{OLTEMODLIN-009}{అమూర్త వాక్యం $!E$ను, అదే
నిరూపణలో సంయోగించే మొదటిస్థాయి వాక్యాలకు కూడా వాడటం వల్ల
బంధిత చర పేరు ఢీకొంటుంది. అధ్యాయం ముందే మొదటిస్థాయి సూత్రాలకు
$!A,!B,\ldots$ను, అమూర్త వాక్యాలకు $!E,!F,\ldots$ను వేరు చేసింది;
కాబట్టి ఇక్కడ బంధిత మొదటిస్థాయి వాక్యానికి $!A$ను వాడాం.}
ఇప్పుడు $!D=\textstyle\bigvee
\Setabs{!D_{\Struct{M}}}{\Struct{M}\models_L!E}$ అని ఉంచండి. ఇక్కడ
ప్రతి తుల్యతా-వర్గం నుంచి ఒక్క ప్రతినిధినే తీసుకున్నందున ఈ వియోగం కూడా
పరిమితమే.
\sourcecorrection{OLTEMODLIN-020}{ఆంగ్ల మూలం “అన్ని” సత్యమైన
పరిమాణక-ర్యాంకు వాక్యాల అనంత సంయోగాన్ని, అన్ని $!E$-నమూనాలపై
వియోగాన్ని రాసి, అవి తార్కిక తుల్యత వరకు పరిమితమని మాత్రమే అంటుంది.
వాస్తవంగా పరిమిత సూత్రాలు రావడానికి తుల్యతా-వర్గాల నుంచి
ప్రతినిధులను ఎంచుకున్న నిర్మాణాన్ని స్పష్టం చేశాం.}

ఈ నిర్మాణంతో $\Mod(L){!E}=\Mod(L){!D}$ వస్తుంది. నిజానికి,
$\Struct{N}\models_L!D$ అయితే, ఏదో $\Struct{M}\models_L!E$కు
$\Struct{N}\models!D_{\Struct{M}}$; అందువల్ల
  లెమ్మా పరికల్పన వల్ల $\Struct{N}\models_L!E$. మరోవైపు
  $\Struct{N}\models_L!E$ అయితే
$!D_{\Struct{N}}$ అనేది $!D$లోని ఒక వియోగకం; అలాగే
$\Struct{N}\models!D_{\Struct{N}}$, కాబట్టి
$\Struct{N}\models_L!D$.
\sourcecorrection{OLTEMODLIN-010}{ఆంగ్ల మూలం ఈ రెండవ దిశలో
$!D_\Struct{N}$ అని రాసి, $\Struct{N}$ను ఉపసూచిక వెలుపల ఉంచుతుంది.
ముందు నిర్వచించిన $!D_{\Struct{M}}$కు అనుగుణంగా రెండు చోట్లా
$!D_{\Struct{N}}$గా సరిచేశాం.}
\end{proof}

\begin{thm}[లిండ్‌స్ట్రోమ్ సిద్ధాంతం]
  \ollabel{thm:lindstrom} $\tuple{L,\models_L}$ ఒక సాధారణ తర్కం; దానికి
  సంహతత్వ, లొవెన్‌హైమ్--స్కోలెమ్ ధర్మాలు ఉన్నాయని అనుకుందాం. అప్పుడు
  $\tuple{L,\models_L}\le\tuple{F,\models}$ (అందువల్ల
  $\tuple{L,\models_L}$ మొదటిస్థాయి తర్కానికి తుల్యం).
  \sourcecorrection{OLTEMODLIN-011}{సాధారణత్వం లేకుండా ఈ
  లక్షణీకరణకు అవసరమైన బూలియన్, విస్తరణ, సమరూపతా, సాపేక్షీకరణ,
  పరిమాణక ధర్మాలు అందుబాటులో ఉండవు; అంతేకాక ముందరి గమనికలో
  ఉపయోగించిన $\tuple{F,\models}\leq\tuple{L,\models_L}$ కూడా
  సాధారణత్వంపైనే ఆధారపడుతుంది. అందువల్ల సిద్ధాంత పరికల్పనలో
  “సాధారణ”ను చేర్చాం.}
\end{thm}

\begin{proof}
\olref{lem:lindstrom} వల్ల, ప్రతి $!E\in L(\Lang{L})$కు, ఏ రెండు
\tetoken{నిర్మాణాలు}{!!{structure}s} $\Struct{M}$, $\Struct{N}$లకైనా
  $\Struct{M}\equiv_n\Struct{N}$ అయినప్పుడు $\Struct{M}$, $\Struct{N}$లు
  $!E$ విషయంలో ఏకీభవించే
$n\in\Nat$ ఉందని చూపడం చాలు. అప్పుడు $!E$ మొదటిస్థాయి
\tetoken{వాక్యానికి}{!!{sentence}} తుల్యం; దాంతో
$\tuple{L,\models_L}\le\tuple{F,\models}$ వస్తుంది. పరిమిత, శుద్ధ
సంబంధాత్మక \tetoken{భాష}{!!{language}}~$\Lang{L}$లో పనిచేస్తున్నందున,
\olref[bas][pis]{thm:b-n-f} ప్రకారం
$\Struct{M}\equiv_n\Struct{N}$ అనే మాటను
$I_n(\emptyseq,\emptyseq)$ అనే బీజగణిత వాక్యంతో మార్చవచ్చు.
\sourcecorrection{OLTEMODLIN-022}{పాక్షిక-సమరూపత అధ్యాయంలోని
నిర్వచనం ఖాళీ అనుక్రమానికి ప్రత్యేక గుర్తు $\emptyseq$ను వాడుతుంది.
ఆంగ్ల మూలం ఇక్కడ రెండు వాదాలుగా సమితి-ఖాళీ $\emptyset$ను రాసింది;
నిర్వచనంతో సరిపడే $\emptyseq$ను పునరుద్ధరించాం.}

ఇప్పుడు $!E$ను స్థిరపరచి, వైరుధ్యం కోసం, ప్రతి $n$కు
\tetoken{నిర్మాణాలు}{!!{structure}s} $\Struct{M}_n$,
$\Struct{N}_n$ ఉండి $I_n(\emptyseq,\emptyseq)$ అయినా,
$\Struct{M}_n\models_L!E$, $\Struct{N}_n\not\models_L!E$ అవుతాయని
అనుకుందాం. సమరూపతా ధర్మం వల్ల అన్ని $\Struct{M}_n$లు భాషలోని
స్థిరసంకేతాలకు ఒకే మూలకాలను అర్థనిర్దేశాలుగా ఇస్తాయని అనుకోవచ్చు.
పరిమిత భాషలో పరమాణు \tetoken{వాక్యాల}{!!{sentence}s} సంఖ్య
పరిమితమే; కాబట్టి \tetoken{వాక్యాల}{!!{sentence}s} విషయంలో అవన్నీ
ఒకే సత్యతా విలువలు ఇచ్చే ఉపక్రమాన్ని ఎంచుకోవచ్చు. సూచికలను మళ్లీ
పేరుపెట్టి, ముందటి $\Struct{M}$ల ఉపక్రమం నుంచి ఎంచుకున్న ఉపక్రమంలో
అసలు ర్యాంకు కనీసం కొత్త సూచికకు
సమానమని కూడా అనుకుందాం; ఎందుకంటే $I_m$ అయితే $m\ge n$కు $I_n$
వస్తుంది. ఇప్పుడు \tetoken{నిర్మాణం}{!!{structure}}~$\Struct{M}$ను
అన్ని $\Struct{M}_n$ల సమ్మేళనంగా తీసుకోండి: ప్రతి $\Struct{M}_n$ను
ఉపనిర్మాణంగా కలిగిన అనన్య కనిష్ఠ నిర్మాణం. స్థిరసంకేతాల సామ్యమైన
అర్థనిర్దేశాలు, పరమాణు వాక్యాలపై ఏకీభవన వల్ల ఈ సమ్మేళనం బాగా
నిర్వచితం. \sourcecorrection{OLTEMODLIN-012}{ఆంగ్ల మూలం
ఉపక్రమం, సంయుక్త మూలకాల సముదాయాల అనుకూలత, వాటి కనిష్ఠ సమ్మేళనాలను
నేరుగా తీసుకుంటుంది. ర్యాంకు తగ్గే ధర్మంతో ఉపక్రమాన్ని మళ్లీ
సూచికీకరించడం, స్థిరసంకేతాల వద్ద మాత్రమే సమానీకరణతో మిగతా
మూలకాలను వియుక్తంగా ట్యాగ్‌చేయడం, పరమాణు వాక్యాలపై ఏకీభవనతో
సంబంధాల సమ్మేళనం బాగా నిర్వచితం కావడం, తరువాతి పరిసర నిర్మాణంలో
మూలకాలు, అనుక్రమాలు, ప్రకృత సంఖ్యల రకాలను పరస్పరం వియుక్తంగా
ట్యాగ్‌చేయడం కూడా స్పష్టం చేశాం.}
\olref[lsp]{thm:abstract-p-isom}లోని నిర్మాణాన్ని అనుసరించి,
$\Struct{M}$ను, $\Domain{M}\cup\Domain{M}^{<\omega}$ అనే
\tetoken{వ్యక్తి క్షేత్రం}{!!{domain}} గల, అనుసంధాన
విధేయ సంకేతాలు $P$, $Q$ను చేర్చిన
\tetoken{భాష}{!!{language}}లోని విస్తరణ $\Struct{M}^*$గా
నిర్మించండి.

అలాగే $\Struct{N}_n$, $\Struct{N}$, $\Struct{N}^*$లను నిర్వచించండి.
ఇప్పుడు \tetoken{నిర్మాణం}{!!{structure}}~$\Struct{K}$ను తీసుకుందాం:
దాని \tetoken{వ్యక్తి క్షేత్రం}{!!{domain}} అనేది $\Struct{M}^*$,
$\Struct{N}^*$ల \tetoken{వ్యక్తి క్షేత్రాల}{!!{domain}s}, అలాగే
ప్రకృత సంఖ్యలు $\Nat$; $\Nat$పై సహజ క్రమం $\le$ ఉంటుంది. దాని
\tetoken{భాష}{!!{language}}లో $\Domain{M}$, $\Domain{N}$,
$\Domain{M}^{<\omega}$, $\Domain{N}^{<\omega}$లను, అలాగే
$\Struct{M}_n$, $\Struct{N}_n$ల క్షేత్రాలను గుర్తించే
\tetoken{వ్యక్తి క్షేత్రాల}{!!{domain}s} సూచకాలు ఉంటాయి; అవి
\begin{align*}
  \Domain{M_n} & = \Setabs{a \in \Domain{M}}{R(a,n)}; &
  \Domain{N_n} & = \Setabs{a \in \Domain{N}}{S(a,n)}; \\
  \Domain{M}^{<\omega}_n & = \Setabs{a \in \Domain{M}^{<\omega}}{R(a,n)}; &
  \Domain{N}^{<\omega}_n & = \Setabs{a \in \Domain{N}^{<\omega}}{S(a,n)}.
\end{align*}
అని నిర్వచించబడతాయి. వియుక్త ట్యాగ్‌లతో అన్ని నాలుగు రకాలు
పరస్పరం కలవకుండా ఎంచుకున్నామని భావించండి.
ఈ \tetoken{నిర్మాణం}{!!{structure}}~$\Struct{K}$కు మరో
త్రిస్థాన సంబంధం $J$ ఉంటుంది; $J(n,\mathbf{a},\mathbf{b})$ అవుతుంది
అప్పుడే $I_n(\mathbf{a},\mathbf{b})$.

\sourcecorrection{OLTEMODLIN-013}{ఆంగ్ల మూలం ఇక్కడ పరిసర
నిర్మాణానికీ $\Struct{M}$ అనే పేరే పెట్టి, తరువాత దాని నమూనాకూ
$\Struct{M}^*$ అనే పేరే వాడుతుంది; ఇది ముందే నిర్మించిన వైపు
$\Struct{M}^*$తో ఢీకొంటుంది. పరిసర నిర్మాణాన్ని $\Struct{K}$గా,
సంహతత్వ నమూనాను $\Struct{K}^*$గా వేరు చేశాం.}
ఇప్పుడు $\Lang{L}$ను $R$, $S$, $J$ మొదలైన సంకేతాలతో విస్తరించిన
భాషలో, $\le$ అనేది మొదటి మూలకం ఉన్నా చివరి మూలకం లేని వివిక్త
రేఖీయ క్రమమని, ప్రతి $n$కు $\Struct{M}_n\models_L!E$,
$\Struct{N}_n\not\models_L!E$ అని, ఇంకా
$J(n,\mathbf{a},\mathbf{b})$ అవుతుంది అప్పుడే
$I_n(\mathbf{a},\mathbf{b})$ అని చెప్పే
\tetoken{ఒక వాక్యం}{!!a{sentence}}~$!D$ ఉంది. ఇది
\tetoken{భాషలో}{!!{language}} నైరూప్య తర్కపు వాక్యంగా
సాపేక్షీకరణతో నిర్మించబడుతుంది.
\sourcecorrection{OLTEMODLIN-014}{ఆంగ్ల మూలం యథేచ్ఛ
అమూర్త $!E$కు $\models$ అనే సాధారణ మొదటిస్థాయి సంతృప్తి సంకేతాన్ని
వాడుతుంది. ఇక్కడ సరైన సంబంధం $\models_L$; దానిని $\Struct{M}_n$
మరియు $\Struct{N}_n$ రెండింటి వద్ద స్పష్టం చేశాం.}

ఇప్పుడు ఈ $!D$తో కలిసి కొత్త స్థిరసంకేతం $d$ను తీసుకొని
\[
\Delta=\{\overline{n}<d:n\in\Nat\}
\]
అని ఉంచండి. ప్రతి పరిమిత ఉపసమితి $!D\cup\Delta_0$కు,
ముందటి ప్రమాణ నమూనా $\Struct{K}$లో $d$ను $\Delta_0$లోని అన్ని
సంఖ్యలకన్నా పెద్దదైన ప్రమాణ సంఖ్యగా అర్థనిర్దేశం చేస్తే నమూనా
ఉంటుంది. కాబట్టి సంహతత్వ ధర్మం వల్ల $!D\cup\Delta$కు
ఒక నమూనా $\Struct{K}^*$ ఉంది; అందులో $d$ విలువను $n^*$గా
రాస్తే, $n^*$ ప్రతి ప్రమాణ సంఖ్యకన్నా పెద్దది, అంటే అప్రామాణికం.
\sourcecorrection{OLTEMODLIN-015}{ఒక్క $!D$కు నమూనా తీసుకుంటే
ప్రమాణ క్రమమే నమూనాగా ఉండగలదు; దాంతో అప్రామాణిక $n^*$
స్వయంగా రావడం లేదు. కొత్త స్థిరసంకేతం $d$తో
$\{\overline n<d:n\in\Nat\}$ రకాన్ని చేర్చాం. ప్రతి పరిమిత భాగం
ప్రమాణ పెద్ద సంఖ్యతో సంతృప్తిపరచబడుతుంది; సంహతత్వం తరువాత లభించే
$d$ విలువే కావలసిన అప్రామాణిక సూచిక.}
అప్పుడు $\Struct{K}^*$లో $\Struct{M_{n^*}}$,
$\Struct{N_{n^*}}$ అనే ఉప\-\tetoken{నిర్మాణాలు}{!!{structure}s} ఉండి
$\Struct{M_{n^*}}\models_L!E$, $\Struct{N_{n^*}}\not\models_L!E$.
కానీ $\Domain{M_{n^*}}$, $\Domain{N_{n^*}}$లలోని $k$-పొడవు
అనుక్రమాల జతకు ఈ $\mathcal{I}$లో $\tuple{\mathbf{a},\mathbf{b}}\in\mathcal{I}$ అవుతుంది
అప్పుడే $J(n^*-k,\mathbf{a},\mathbf{b})$ అని ఉంచండి; ఇక్కడ $k$ అనేది
$\mathbf{a}$, $\mathbf{b}$ల పొడవు. ప్రతి ప్రమాణ $k$కు
$n^*-k>0$ కాబట్టి
$\mathcal{I}$ పాక్షిక సమరూపతా ముందుకు-వెనుకకు ధర్మాన్ని
నెరవేరుస్తుంది; అందువల్ల
$\Struct{M_{n^*}}\simeq_p\Struct{N_{n^*}}$. కానీ
\olref[lsp]{thm:abstract-p-isom} ప్రకారం $\Struct{M_{n^*}}$,
$\Struct{N_{n^*}}$లు $L$-వాక్య తుల్యాలు,
ఇది వైరుధ్యం.
\end{proof}

\end{document}
